\subsection{Le criticità energetiche dell'esecuzione all'edge}

Nonostante i benefici in termini di latenza e prossimità
ai dati, l'esecuzione all'edge introduce criticità significative.
I nodi edge sono tipicamente caratterizzati da risorse
computazionali limitate, eterogeneità hardware e condizioni
operative variabili. A differenza dei data center cloud,
progettati per offrire elevata capacità e ridondanza, i nodi
edge devono operare con vincoli stringenti su CPU, memoria e,
soprattutto, energia disponibile.

In molti scenari applicativi -- dispositivi alimentati a batteria,
infrastrutture soggette a limiti di potenza, contesti IoT -- l'energia
non rappresenta soltanto un costo operativo, ma una risorsa limitata
e dinamica. La letteratura sul resource scheduling in ambienti edge
evidenzia come la gestione efficiente delle risorse, inclusa
l'energia, costituisca una delle principali sfide ancora aperte
\cite{Luo2021_EdgeSchedulingSurvey,Kar2023_OffloadingFederatedSurvey}.

In tali contesti non è sufficiente adottare politiche che
minimizzano il consumo medio di energia: diventa necessario
controllare l'energia associata alle singole esecuzioni,
garantendo il rispetto di budget energetici prefissati
e potenzialmente variabili nel tempo.
